\documentclass[10pt,a4paper]{amsart}
\usepackage[margin=19mm]{geometry}
\usepackage{amsmath,amssymb,mathtools}
\usepackage[scr=rsfs,cal=euler]{mathalfa}
\usepackage{xfrac}
\usepackage[unicode,hidelinks,bookmarks=false]{hyperref}
\newcommand{\ie}{i.e.\ }
\newcommand{\cf}{cf.\ }
\newcommand{\resp}{respectively\ }
\newcommand{\Subsection}{\subsection}
\providecommand{\nobreakdash}{\nobreak\hbox{-}}
\setlength{\emergencystretch}{2em}
\multlinegap0pt
\usepackage{bm}
\usepackage{bbm}
\usepackage{dsfont}
\usepackage{xspace}
\usepackage{cancel}
\usepackage{mathtools}
\usepackage{amsthm}
\makeatletter
\@ifundefined{theorem}{\newtheorem{theorem}{Theorem}}{}
\@ifundefined{proposition}{\newtheorem{proposition}[theorem]{Proposition}}{}
\makeatother
\title[On Semisymmetric Height and a Multidimensional Generalizatio]{On Semisymmetric Height and a Multidimensional Generalization of Weighted Catalan Numbers}
\author{Ryota Inagaki, Dimana Pramatarova}
\date{Source: arXiv:2604.04900v1 (CC BY 4.0)}
\begin{document}
\maketitle

\noindent\emph{Source and licence.} On Semisymmetric Height and a Multidimensional Generalization of Weighted Catalan Numbers. Version arXiv:2604.04900v1 (CC BY 4.0), distributed under the Creative Commons Attribution 4.0 International licence, as recorded in the arXiv licence metadata for this exact version and shown on the abstract page https://arxiv.org/abs/2604.04900v1. Full rights notice: licenses/arxiv-2604.04900-CC-BY-4.0.txt.\par\smallskip

\noindent\textit{Editorial scope.} Counted unit: the Proposition whose source label is \texttt{prop:NotTightExample} in the article (source0.tex lines 800--802, source label \texttt{prop:NotTightExample}), delivered together with the article's own complete original proof of it (source0.tex lines 803--824, one \texttt{proof} environment of 22 lines). No mathematics is written, repaired or re-derived: statement and proof are the article's own text line for line. Required context retained uncounted: none. Every definition, display and label of those lines is retained unchanged. Omitted from this package and not redistributed: 1--799, 825--984. Nothing inside a retained excerpt is omitted or altered: every retained body is delivered exactly as the article's lines are written, so no omission receipt of any kind accompanies this package. Citations: the retained excerpts carry no citation command at all, so no bibliography entry of the article is transcribed and no reference is redistributed. Reference adaptation: none was needed and none was made; every reference inside the delivered unit resolves inside it, so no unresolved reference or citation is produced. Printed slips kept literal: no printed word, symbol, hypothesis or number of the counted statement or of its proof was altered. Layout and class: the article's own class is \texttt{article} with packages including amsmath, amssymb, amscd, amsthm, amsfonts, tikz, graphicx, hyperref, float, comment, caption, subcaption, ytableau; the wrapper is the lane's pinned \texttt{amsart} editorial wrapper at 10pt on A4, which restates the theorem-like environments the retained text needs and adds the article's own macro definitions that the retained text needs (none), with extra wrapper packages (bm, bbm, dsfont, xspace, cancel, mathtools). The delivered numbering is the unit's own. Assets: no retained span contains a figure environment, a graphics-inclusion command, a PDF-page inclusion, a table or a TikZ picture; no image file is copied into this package, the package declares no asset dependency and no image file is redistributed with it. Licence: version arXiv:2604.04900v1 of this article is distributed under the Creative Commons Attribution 4.0 International licence, as recorded in the arXiv licence metadata for this exact version and shown on the abstract page https://arxiv.org/abs/2604.04900v1. A full rights notice is delivered at licenses/arxiv-2604.04900-CC-BY-4.0.txt. Characters: every retained excerpt is pure ASCII, so the delivered engine, XeLaTeX with its default Latin Modern text font, reports no missing character.
\begin{proposition}\label{prop:NotTightExample}
    For each $n \geq 1$, we have $D'_{3, 2n, n} = \binom{3n}{n} - 2\binom{3n}{n-1} + \binom{3n}{n-2}$.
\end{proposition}

\begin{proof}
     Observe that the expression $\binom{3n}{n} - 2\binom{3n}{n-1} + \binom{3n}{n-2}$ is the number of paths from $(0, 0)$ to $(0, 0)$ within $\mathbb{Z}_{\geq 0}^2$ consisting of $3n$ steps from $\{A=(1, 0), B= (-1, 1), C=(0, -1)\}$. Call the set of such paths $Q_{n}$. Thus, it suffices to show that the set of $3$-dimensional balanced ballot paths counted by $D'_{3, 2n, n}$, i.e., the $3n$ steps with exact semisymmetric height  $2n$, is equinumerous to $Q_n$. We prove this fact via a bijective proof.

    Consider the following map $f$ from the set  of $3$-dimensional balanced ballot paths counted by $D'_{3, 2n, n}$ to the set $Q_n$: for balanced ballot path $P$ in the domain, we construct path $q \in Q_n$ by letting the $i$th step of $q$ be $A$ if the $i$th step of $P$ is $\vec{e}_1$, $B$ if the $i$th step of $P$ is $\vec{e}_2$, and $C$ if the $i$th step of $P$ is $\vec{e}_3$. 
    
    We confirm that, for any 3-dimensional balanced ballot path $P$ of length $3n$ with semisymmetric height $2n$, the path  $f(P)$ is in $Q_n$. First, the intermediate points of $Q_n$ are all in $\mathbb{Z}_{\geq 0}^2$. Indeed, ballot property of $P$ guarantees that, in any prefix of the sequence $f(P)$, the number of steps equal to $A = (1, 0)$ is at least that of steps equal to $B = (-1, 1)$, and that the number of steps equal to $B = (-1, 1)$ is at least that of steps equal to $C = (0, 1)$.
    
    
    Thus, to show $f(P)$ is in $Q_n$, it suffices to show that the last step equal to $A$ in path $f(P)$ is before the first step equal to $C$. We prove an equivalent statement: the first $\vec{e}_3$ step (i.e., the first semisymmetric down-step) in $P$ must be after the last $\vec{e}_1$ step (i.e., last semisymmetric down-step). Suppose, for the sake of contradiction, that this statement is not true. Let $i_{max} \in \{0, 1, \dots, kn\}$ be the smallest possible integer so that the semisymmetric height of $P$, i.e., the maximum possible semisymmetric height attainable by an intermediate point of $P$, is that of the $i_{max}$th intermediate point of $P$. Let $\vec{v}_{i_{max}}=(x_1, x_2, x_3)$ be the $i_{max}$th intermediate point of $P$. We know that $i_{max}$th step must not be a semisymmetric down-step; otherwise, the $(i_{max}-1)$th intermediate point of $P$ must have a higher semisymmetric height than the $i_{max}$th. We also know the $i_{max}$th step of $P$ must be an up-step; otherwise the $(i_{max}-1)$th intermediate point would have a semisymmetric height equal to that of $P$, contradicting the definition of $i_{max}$. We then obtain that $g_3(\vec{v}_{i_{max}})=2x_1-2x_3 < 2n$. If the $i_{max}$th step is not the last semisymmetric up-step, then we find that $g_3(\vec{v}_{i_{max}})=2x_1-2x_3 \leq 2x_1 < 2n$. Otherwise, if the $i_{max}$th step is the last semisymmetric up-step, we see $g_3(\vec{v}_{i_{max}})=2x_1-2x_3 < 2x_1 = 2n$, because of the semisymmetric down-step that came before the last semisymmetric up-step. Therefore, $g_k(\vec{v}_{i_{max}}) < 2n$ is the semisymmetric height of $P$, and a contradiction is reached. 

    In sum, we have that $f(P)$ is in $\mathbb{Z}_{\geq 0}^2$, and the last step equal to $A$ in $f(P)$ is before the first equal to $C$.  Hence, the path $f(P)$ is in $Q_n$.
    
    We now show that the function $f$ is a bijection because it has an inverse function $g$. Consider the function $g$, defined as follows: for $q \in Q_n$, if the $i$th step of $q$ is $A$ then the $i$th step of $g(q)$ is $\vec{e}_1$, if the $i$th step of $q$ is $B$ then the $i$th step of $g(q)$ is $\vec{e}_2$, and if the $i$th step of $q$ is $C$ then the $i$th step of $g(q)$ is $\vec{e}_3$. 
    
    We observe that the range of $g$ is a subset of $Q_n$. Firstly, there are $n$ steps in $g(q)$ equal to $\vec{e}_1, \vec{e}_2$, and $\vec{e}_3$ since there are $n$ steps equal to $A$, $B$, and $C$ in $q$. Furthermore, path $g(q)$ exhibits the ballot property. Otherwise, for some prefix of $q$, there are 1) more instances of $B$  than $A$ or 2) more of $C$ than $A$, and that would imply that $q$ is not in $\mathbb{Z}_{\geq 0}^2$, a contradiction of how $q \in Q_n$.
    
    Furthermore, the balanced ballot path $g(q)$ has a semisymmetric height of exactly $2n$. Indeed, the last step equal to $A$ in path $q$ is before the first step equal to $C$, which means the last instance of $\vec{e}_1$ in $g(q)$ is before the first of $\vec{e}_3$. Consequently, the semisymmetric height of the intermediate point resulting from the last $\vec{e}_1$ step is $2n$.

    Finally, function $g$ is the inverse of $f$. Indeed, for any $q \in Q_n$ and  3-dimensional balanced ballot path $P$ of $3n$ steps with semisymmetric height $2n$, we obtain $f(g(q)) = q$ and $g(f(P))= P$ because $f$ and $g$ relabel steps.

    Because $f$ is a bijection from the set of $3$-dimensional balanced ballot paths of length $3n$ with semisymmetric height $2n$ to $Q_n$, we find that $D_{3, 2n, n}' = |Q_n| = \binom{3n}{n} - 2\binom{3n}{n-1}+ \binom{3n}{n-2}$.
\end{proof}

\end{document}
